\section*{Chapter \ref{ch:b3:electronic-spectroscopy} --- Electronic Spectroscopy and Photophysics}

\begin{solution}{exo:b3:electronic-spectroscopy:1}
(a) Forbidden by the spin rule ($\Delta S = 1$). (b) Forbidden by the \omterm{thm:b3:electronic-spectroscopy:laporte}{Laporte rule}
($g \to g$); seen weakly through vibronic coupling. (c) Allowed: $B_{1u}$ is the
representation of $z$ in $D_{2h}$. (d) Forbidden by symmetry: $A_2$ is none of $x$,
$y$, $z$ in $C_{2v}$; the band is weak.
\end{solution}

\begin{solution}{exo:b3:electronic-spectroscopy:2}
$10^7/349 - 10^7/450 = 28653 - 22222 = \qty{6431}{cm^{-1}}$, \qty{0.80}{eV}.
\end{solution}

\begin{solution}{exo:b3:electronic-spectroscopy:3}
$A = 5700 \times 1 \times 1.0\times10^{-5} = 0.057$; absorbed fraction $1 - 10^{-0.057} =
0.12$.
\end{solution}

\begin{solution}{exo:b3:electronic-spectroscopy:4}
$k_r = \Phi_F/\tau = \qty{1.5e8}{s^{-1}}$, $\tau_r = \qty{6.7}{ns}$; $k_{nr} = (1 - \Phi_F)/\tau
= \qty{1.0e8}{s^{-1}}$.
\end{solution}

\begin{solution}{exo:b3:electronic-spectroscopy:5}
$\int\varepsilon\,\dd\tilde\nu = 1.0645 \times 5700 \times 5000 = \num{3.03e7}$, so $f \approx
4.32\times10^{-9} \times 3.03\times10^7 = 0.13$: an allowed transition of moderate
strength.
\end{solution}

\begin{solution}{exo:b3:electronic-spectroscopy:6}
The Poisson factor $\eu^{-S}S^v/v!$ is maximal at $v = \lfloor S\rfloor$. $S = 0.3$: the
0--0 line itself. $S = 1.5$: $v' = 1$, 1.5 times the 0--0 line. $S = 5$: $v' = 4$
and 5 equally, 26 times the 0--0 line.
\end{solution}

\begin{solution}{exo:b3:electronic-spectroscopy:7}
$\sum k = \qty{4.2e8}{s^{-1}}$: $\tau = \qty{2.4}{ns}$, $\Phi_F = 0.24$, $\Phi_{\mathrm{ISC}} = 0.71$
(and $\Phi_{\mathrm{IC}} = 0.05$).
\end{solution}

\begin{solution}{exo:b3:electronic-spectroscopy:8}
$\Phi = 0.546 \times 0.46 \times (0.050/0.040) \times (1.4266/1.333)^2 = 0.36$.
\end{solution}

\begin{solution}{exo:b3:electronic-spectroscopy:9}
$\tau_0/\tau = 1.429$ and 1.860: linear in $[Q]$, so the lifetime itself is
shortened: \omterm{def:b3:electronic-spectroscopy:quencher}{dynamic quenching}. $K_{SV} = 43$ L/mol (both points), $k_q = 43/8.0
\times 10^{-9} = \qty{5.4e9}{L.mol^{-1}.s^{-1}}$.
\end{solution}

\begin{solution}{exo:b3:electronic-spectroscopy:10}
\omterm{def:b3:electronic-spectroscopy:quencher}{Static quenching}: the molecules that still emit live as long as before; half of
them are bound, in the ground state, in a non-fluorescent complex. With an
association constant $K_a$, the free fraction is $1/(1 + K_a[Q])$ and $I_0/I = 1 +
K_a[Q]$, the same form as Stern--Volmer but with $\tau$ unchanged.
\end{solution}

\begin{solution}{exo:b3:electronic-spectroscopy:11}
$\tau_r = \tau/\Phi_F = \qty{14}{ns}$; $\Phi_T = 1 - 0.36 = 0.64$. The measured lifetime is
shortened by the competing \omterm{def:b3:electronic-spectroscopy:radiationless}{intersystem crossing}; $\tau_r$ is what emission alone
would give.
\end{solution}

\begin{solution}{exo:b3:electronic-spectroscopy:12}
$\frac12\langle\alpha\beta - \beta\alpha|\alpha\beta + \beta\alpha\rangle = \frac12(1 + 0 - 0 - 1) = 0$, using
$\langle\alpha|\alpha\rangle = \langle\beta|\beta\rangle = 1$, $\langle\alpha|\beta\rangle = 0$ for each electron.
The dipole \omterm{def:b3:quantum-model-systems:operator}{operator} does not touch spin, so the transition moment contains this
zero factor: $T_1 \leftarrow S_0$ has no intensity unless \omterm{def:b3:many-electron-atoms:spin-orbit}{spin--orbit coupling} mixes
the states.
\end{solution}

\begin{solution}{pb:b3:electronic-spectroscopy:1}
\textbf{1.} $A = 5700 \times 1 \times 1.0\times10^{-4} = 0.57$.
\textbf{2.} $1 - 10^{-0.57} = 0.73$.
\textbf{3.} $\varepsilon$ of a few thousand: an allowed $\pi\to\pi^*$ transition, of moderate
strength.
\textbf{4.} $1239.8/349 = \qty{3.55}{eV}$.
\textbf{5.} It absorbs in the ultraviolet only; the visible light passes.
\textbf{6.} To keep the absorbed light proportional to the concentration and avoid
re-absorption of the emitted light (inner-filter effects).
\textbf{7.} $S_0 \to S_1$ (absorption, then \omterm{def:b3:electronic-spectroscopy:radiationless}{vibrational relaxation}); from $S_1$:
\omterm{def:b3:electronic-spectroscopy:luminescence}{fluorescence}, \omterm{def:b3:electronic-spectroscopy:radiationless}{internal conversion}, \omterm{def:b3:electronic-spectroscopy:radiationless}{intersystem crossing} to $T_1$.
\textbf{8.} $k_r = 0.546/19 \times 10^{-9} = \qty{2.87e7}{s^{-1}}$; $\tau_r = \qty{35}{ns}$.
\textbf{9.} $(1 - 0.546)/\tau_0 = \qty{2.39e7}{s^{-1}}$.
\textbf{10.} $28653 - 22222 = \qty{6431}{cm^{-1}}$.
\textbf{11.} Emission starts from the relaxed $S_1$ and ends on vibrationally excited
levels of $S_0$, after the molecule and solvent have relaxed: the photon has less
energy, here enough to fall in the visible.
\textbf{12.} $\Phi_F \times (349/450) = 0.42$: 42\,\% of the absorbed energy.
\textbf{13.} The points rise by 0.594 per \qty{0.010}{mol/L}: a straight line through 1.
\textbf{14.} Least squares through the five points: slope $K_{SV} = \qty{59.4}{L/mol}$,
intercept 1.000.
\textbf{15.} $k_q = K_{SV}/\tau_0 = \qty{3.1e9}{L.mol^{-1}.s^{-1}}$.
\textbf{16.} $\tau = \tau_0/(1 + 59.4 \times 0.040) = \qty{5.6}{ns}$.
\textbf{17.} Measure lifetimes: for \omterm{def:b3:electronic-spectroscopy:quencher}{dynamic quenching} $\tau_0/\tau$ falls on the same line
as $I_0/I$.
\textbf{18.} Tonic water contains no chloride to quench it, and is acidic enough to
keep quinine protonated, its fluorescent form.
\textbf{19.} $k_d = 8 \times 8.314 \times 298/(3 \times 0.890\times10^{-3}) =
\qty{7.4e6}{m^3.mol^{-1}.s^{-1}} = \qty{7.4e9}{L.mol^{-1}.s^{-1}}$.
\textbf{20.} $k_q$ is a little less than half of $k_d$.
\textbf{21.} About 42\,\%.
\textbf{22.} Slower: $k_d \propto 1/\eta$, and $k_q$ can only fall with it.
\textbf{23.} \textbf{$k_q/k_d = 0.42$}: chloride quenches quinine at almost half the
rate of diffusion-controlled encounters.
\end{solution}
